\documentclass[border=4pt]{standalone}
\usepackage{amsmath,amssymb,mathtools,xcolor,varwidth}
\definecolor{arcblue}{HTML}{1E6FE0}
\begin{document}
\begin{varwidth}{28em}
\large\bfseries
We resume the very figure of the last lemma. From the point $A$ springs \textcolor{arcblue}{the arc $ACB$},      fixed in position, with smooth continuous curvature; across it runs the chord $AB$, and grazing      the curve at $A$ stretches the tangent $AD$. As the point $B$ slides toward $A$, all three of      these --- \textcolor{arcblue}{arc, chord, and tangent} --- shrink together toward nothing, and the question is:      in what \emph{proportion} do they vanish? Newton's claim is bold and exact: their ultimate ratio,      any one to any other, is the \textcolor{arcblue}{ratio of equality} --- $1:1:1$. The trouble is that three      quantities all collapsing to zero cannot be compared directly, so we need a clever device to hold      their proportions still while we examine them.
\end{varwidth}
\end{document}
